\documentclass[11pt]{amsart}
\usepackage{amssymb}
\usepackage{url}
\usepackage[margin=1.2in]{geometry}

\newtheorem{theorem}{Theorem}
\newtheorem{proposition}[theorem]{Proposition}
\newtheorem{lemma}[theorem]{Lemma}
\newtheorem{corollary}[theorem]{Corollary}
\theoremstyle{remark}
\newtheorem{remark}[theorem]{Remark}

\title{Rank-one Horrocks over a local ring by division}
\author{Claude (Anthropic, Opus 5.5)}
\date{}

% stamp on every page: file name and date/time of compilation
\newcount\stamphour \newcount\stampmin
\stamphour=\time \divide\stamphour by 60
\stampmin=\time \multiply\stamphour by 60 \advance\stampmin by -\stamphour \divide\stamphour by 60
\edef\stamptime{\the\year-\ifnum\month<10 0\fi\the\month-\ifnum\day<10 0\fi\the\day\space
 \ifnum\stamphour<10 0\fi\the\stamphour:\ifnum\stampmin<10 0\fi\the\stampmin}
\AddToHook{shipout/foreground}{\put(0.5\paperwidth,-\dimexpr\paperheight-0.45in\relax)%
 {\makebox[0pt]{\scriptsize\ttfamily\jobname.pdf\quad\stamptime}}}
\begin{document}

\begin{abstract}
Let $A$ be a local ring and $L$ an ideal of $A[X]$ containing a monic
polynomial. We show that if the inclusion $L\subseteq A[X]$ has an
$A$-linear retraction, equivalently if $A[X]/L$ is a projective $A$-module,
then $L$ is generated by a unique monic polynomial. An invertible ideal
containing a monic polynomial has such a retraction, built from division by
a monic polynomial; this gives Horrocks' theorem in rank one. The proof is
an induction on the degree of a monic $F\in L$, with the retraction fixed.
Its core is a lemma on $A$-direct summands of $A[X]/(F)$ stable under
multiplication by $X$: such a summand is $0$ or contains a monic polynomial
of degree $\deg F-1$. The argument is Swan's proof of Horrocks' theorem,
in Lindel's matrix form, specialised to rank one, together with the
Nakayama step of Nashier and Nichols. What is new is that it is
constructive: there is no residue field, Nakayama's lemma is used without
determinants and without the radical, and neither freeness nor Schanuel's
lemma is needed.
\end{abstract}

\maketitle

\section{Statement}

All rings are commutative. $A$ is a local ring: if $a+b$ is a unit, then $a$
or $b$ is a unit. We do not require $0\neq1$: the zero ring satisfies this
condition, and the theorems hold for it trivially. Locality is used only
once, in Lemma~\ref{lem:local}. For $p\in A[X]$ write $p_k$ for the
coefficient of $X^k$, and $A[X]_{<d}\cong A^d$ for the polynomials of degree
$<d$.

\begin{theorem}\label{thm:main}
Let $L$ be an ideal of $A[X]$ containing a monic polynomial, and let
$\rho:A[X]\to L$ be $A$-linear with $\rho(p)=p$ for $p\in L$. Then $L$ is
generated by a monic polynomial, and this polynomial is unique.
\end{theorem}

The retraction $\rho$ exists exactly when $L$ is an $A$-direct summand of
$A[X]$, that is, when $A[X]/L$ is a projective $A$-module; it is then
finitely generated, as a quotient of $A[X]/(F)\cong A^d$ for a monic $F\in L$
of degree $d$. There is a short argument with determinants: $A[X]/L$ is
finitely generated projective over the local ring $A$, hence free, say of
rank $e$ (see~\cite{LQ}); the characteristic polynomial $G$ of
multiplication by $X$ on $A[X]/L$ is monic of degree $e$ and lies in $L$ by
Cayley--Hamilton; and $A[X]/(G)\to A[X]/L$ is a surjection between free
modules of rank $e$, hence an isomorphism, so $L=(G)$. This argument is also
constructive. The proof below uses neither freeness nor determinants.

The rank-one case is Corollary~\ref{cor:rankone}: an invertible ideal
containing a monic polynomial has a retraction. This is Horrocks'
theorem~\cite{Hor} in rank one. It was proved by Bass and Murthy~\cite{BM}
for any commutative ring, not only a local one; see the notes to Chapter~IV
of~\cite{Lam}. In~\cite[Ch.~IV, \S5]{Lam} it is Prop.~IV.5.2, due to
Nashier and Nichols~\cite{NN}, and the start of the induction on the rank
in their proof of Horrocks' theorem (Thm.~IV.5.4). A constructive treatment
of Horrocks' theorem is in Lombardi--Quitt\'e~\cite{LQ}.

The proof below is not a new argument. Specialised to rank one, it is
Swan's proof of Horrocks' theorem via Towber's presentation~\cite[\S IV.3]{Lam78},
\cite[\S IV.3]{Lam}, following~\cite{Fer}, and Lindel's matrix form of
it~\cite{Lin}. Swan's module $M=D\cap P$, with $D$ the polynomials of degree
$<\deg f$, is our $M_F$ (Section~\ref{sec:induction}), the action $a$ of $t$
on it is our $T$, and his map $\beta$, defined by $(t-a)m=f\beta(m)$, is the
top coefficient $w\mapsto w_{d-1}$. Swan's key claim that $\beta$ is nonzero
modulo the maximal ideal becomes $1\in J$, and splitting off $Rx$ with
$\beta(x)$ a unit becomes the passage to a monic polynomial of degree $d-1$.
The Nakayama step on $M_F\cong L/(F)$ is that of Nashier and Nichols. The
differences are these: Swan proves that $\beta$ is nonzero modulo the maximal
ideal by an argument on the residue field~\cite[IV.3.4]{Lam78}, and
Nashier--Nichols use that $k[t]$ is a principal ideal domain and Nakayama's
lemma for the radical; here the residue field never appears, and
Lemma~\ref{lem:coef} together with the determinant-free NAK and locality
replace these steps. Swan obtains the projectivity of $M$ from Schanuel's
lemma and uses that it is free; here an explicit retraction suffices, and
no basis of $M$ is chosen.

The proof has three parts: stable summands give a degree drop
(Section~\ref{sec:stable}); with $\rho$ fixed, induction on the degree gives
the generator (Section~\ref{sec:induction}); for an invertible ideal,
division constructs $\rho$ (Section~\ref{sec:rankone}).

\section{Stable summands}\label{sec:stable}

Fix a monic $F$ of degree $d>0$ and define $T:A[X]_{<d}\to A[X]_{<d}$ by
\[
T(w)=Xw-w_{d-1}F;
\]
the right-hand side has degree $<d$ since $F$ is monic. Identifying
$A[X]_{<d}$ with $A[X]/(F)$ by remainders, $T$ is multiplication by $X$.
Division by $F$ itself is needed only from Section~\ref{sec:induction} on.

In this section $M\subseteq A[X]_{<d}$ is an $A$-submodule with an
$A$-linear map $\pi:A[X]_{<d}\to M$ such that $\pi(w)=w$ for $w\in M$; that
is, $M$ is an $A$-direct summand of $A[X]_{<d}\cong A^d$. Put
$u^{(t)}=\pi(X^t)\in M$ for $t<d$; the superscript indexes these elements,
while subscripts always denote coefficients.

\begin{lemma}[coordinates]\label{lem:span}
Every $w\in M$ is $w=\sum_{t<d}w_t\,u^{(t)}$.
\end{lemma}

\begin{proof}
$w=\pi(w)=\pi(\sum_tw_tX^t)=\sum_tw_t\,\pi(X^t)$, by linearity of $\pi$.
\end{proof}

The coefficients in this expansion are the coefficients of $w$ itself,
which is more than finite generation of $M$. Let $J\subseteq A$ be the ideal
generated by the top coefficients $u^{(t)}_{d-1}$, $t<d$.

\begin{lemma}[coefficients]\label{lem:coef}
If $T(M)\subseteq M$, then $M\subseteq J[X]$.
\end{lemma}

\begin{proof}
By Lemma~\ref{lem:span}, $w_{d-1}=\sum_tw_tu^{(t)}_{d-1}\in J$ for every
$w\in M$. We show $w_k\in J$ for all $w\in M$ by descending induction on
$k$. Let $b=T(w)=Xw-w_{d-1}F\in M$. Comparing coefficients of $X^{k+1}$
gives $w_k=b_{k+1}+w_{d-1}F_{k+1}$, which lies in $J$ since $b_{k+1}\in J$
by induction.
\end{proof}

Lemma~\ref{lem:coef} replaces the step $I\cap\mathfrak m[t]=\mathfrak mI$ in
the proof of \cite[Prop.~IV.5.2]{Lam}, and Swan's argument that $\beta$ is
nonzero modulo the maximal ideal~\cite[IV.3.4]{Lam78}.

By Lemmas~\ref{lem:span} and~\ref{lem:coef}, $u^{(s)}=\sum_tu^{(s)}_t\,u^{(t)}$
with all $u^{(s)}_t\in J$: the module spanned by the $u^{(t)}$, which is
$M$, satisfies $M=JM$.

\begin{lemma}[NAK]\label{lem:nak}
Let $B$ be any commutative ring, $J$ an ideal of $B$, $V$ a $B$-module and
$v_1,\dots,v_m\in V$ with $v_s=\sum_ta_{st}v_t$, all $a_{st}\in J$. Then some
$\mu$ with $\mu-1\in J$ kills every $v_s$.
\end{lemma}

\begin{proof}
We prove more generally: if $\tau-1\in J$ and $\tau v_s=\sum_ta_{st}v_t$ with
$a_{st}\in J$, then some $\mu$ with $\mu-1\in J$ kills every $v_s$. Induction
on $m$; for $m=0$ take $\mu=1$. Let $\sigma=\tau-a_{mm}$, so $\sigma-1\in J$
and $\sigma v_m=\sum_{t<m}a_{mt}v_t$. For $s<m$,
\[
\sigma\tau v_s=\sigma a_{sm}v_m+\textstyle\sum_{t<m}\sigma a_{st}v_t
=\sum_{t<m}(\sigma a_{st}+a_{sm}a_{mt})v_t,
\]
with coefficients in $J$ and $\sigma\tau-1\in J$. By induction some $\mu'$
with $\mu'-1\in J$ kills $v_1,\dots,v_{m-1}$; then $\mu=\mu'\sigma$ kills
$v_m$ too, and $\mu-1=\mu'(\sigma-1)+(\mu'-1)\in J$.
\end{proof}

This is the first half of Theorem~2.2 (NAK) of Matsumura~\cite{Mat}: if
$M=JM$ with $M$ finite, some $\mu\equiv1\bmod J$ kills $M$. As Matsumura
remarks there, it can be proved without determinants, by induction on the
number of generators, as above; he also recalls that Nakayama attributed the
result to Krull and Azumaya. No decision on units is needed. We do not use
the consequence that $M=0$ when $J$ lies in the radical, since $J$ is not
assumed to be in the radical. Locality gives an alternative instead.

\begin{lemma}[local NAK]\label{lem:local}
In Lemma~\ref{lem:nak}, if $B$ is local, then all $v_s$ vanish or $1\in J$.
\end{lemma}

\begin{proof}
Since $\mu+(1-\mu)=1$, $\mu$ or $1-\mu$ is a unit. In the first case $\mu$
kills nothing but $0$. In the second case $1-\mu\in J$ is a unit, so $1\in J$.
\end{proof}

\begin{proposition}[stable summands]\label{prop:stable}
Let $T(M)\subseteq M$. Then $M=0$, or $M$ contains a monic
polynomial of degree $d-1$.
\end{proposition}

\begin{proof}
By Lemmas~\ref{lem:coef} and~\ref{lem:local}, applied to the $u^{(t)}$, all
$u^{(t)}$ vanish or $1\in J$. In the first case $M=0$ by
Lemma~\ref{lem:span}. In the second, write $1=\sum_t\gamma_tu^{(t)}_{d-1}$.
Then $W=\sum_t\gamma_tu^{(t)}=\pi(\sum_t\gamma_tX^t)$ lies in $M$, has degree
$<d$ and top coefficient $W_{d-1}=1$.
\end{proof}

Under the identification of $A[X]_{<d}$ with $A[X]/(F)$, a $T$-stable $M$
is an $A[X]$-submodule of $A[X]/(F)$, that is, $L/(F)$ for an ideal
$L\supseteq(F)$. This is how Proposition~\ref{prop:stable} is used in
Section~\ref{sec:induction}.

\section{The induction}\label{sec:induction}

Let $L$ and $\rho$ be as in Theorem~\ref{thm:main}, and let $F\in L$ be monic
of degree $d$. Write $p=Fq(p)+r(p)$ with $\deg r(p)<d$. Since $F$ is monic,
this decomposition is unique (also for $d=0$), so the quotient $q$ and the
remainder $r$ are $A$-linear. For $\deg w<d$ we have $Xw=F\,w_{d-1}+T(w)$,
so $T(w)=r(Xw)$. Put
\[
M_F=\{w\in L\mid \deg w<d\},\qquad \pi_F=r\circ\rho \text{ on } A[X]_{<d}.
\]
If $p\in L$ then $r(p)=p-Fq(p)\in M_F$. Hence $\pi_F$ maps into $M_F$,
and $\pi_F(w)=r(w)=w$ for $w\in M_F$. Since $L$ is an ideal, $Xw\in L$ for
$w\in M_F$, so $T(w)=r(Xw)\in M_F$. Only division changes with $F$; the
retraction $\rho$ does not.

\begin{lemma}[step]\label{lem:step}
Either $F$ divides every element of $L$, or $L$ contains a monic polynomial
of degree $d-1$.
\end{lemma}

\begin{proof}
If $M_F=0$, then $r(p)=0$ for $p\in L$, so $p=Fq(p)$. If $d=0$, then
$M_F=0$. If $d>0$, Proposition~\ref{prop:stable} applies to $M_F$ and
$\pi_F$.
\end{proof}

\begin{lemma}[uniqueness]\label{lem:unique}
Over any commutative ring, monic polynomials $G$, $H$ with $H=Gb$ and
$G=Ha$ are equal.
\end{lemma}

\begin{proof}
Let $n=\deg G$ and $m=\deg H$. If $m<n$, then $H=G\cdot b+0=G\cdot 0+H$
are two divisions of $H$ by $G$, so $H=0$ by uniqueness of the remainder,
and $G=Ha=0$. The case $n<m$ is symmetric. If $m=n$, then $H-G$ has degree
$<n$, and $H=G\cdot 1+(H-G)=G\cdot b+0$ gives $H-G=0$.
\end{proof}

No comparison of $a$ or $b$ with units is needed; this matters when $A$ has
nilpotents, since then $A[X]$ has units of positive degree. Neither
locality nor the retraction is used here.

\begin{proof}[Proof of Theorem~\ref{thm:main}]
Existence is by induction on the degree $d$ of a monic polynomial $F\in L$,
for all such $F$ at once. By Lemma~\ref{lem:step}, either $F$ generates $L$,
or $L$ contains a monic polynomial of degree $d-1$ and the induction
hypothesis applies. If $G$ and $H$ are monic generators of $L$, each
divides the other, so $G=H$ by Lemma~\ref{lem:unique}.
\end{proof}

Locality only gives a mere disjunction, so the induction gives the mere
existence of a generator. Since monic generators are unique, the type of
monic generators of $L$ is a proposition, and mere existence gives an actual
generator.

\begin{remark}[arbitrary rings]\label{rem:dynamic}
Locality is used only in Lemma~\ref{lem:local}, on $\mu+(1-\mu)=1$; the
division, the retraction, Lemmas~\ref{lem:span} and~\ref{lem:coef} and
Lemma~\ref{lem:nak} hold over any commutative ring. Over an arbitrary $A$,
one can split there instead of deciding, in the style of dynamical
evaluation~\cite[XIV.1]{LQ}: $\mu$ and $1-\mu$ are comaximal, and $L$,
$\rho$ and $F$ localise. Over $A[1/\mu]$, $\mu$ kills $M_F$, so $r(p)=0$ for
$p\in L$ and $F$ generates $L$. Over $A[1/(1-\mu)]$, the element $\mu-1\in J$
is a unit, so $1\in J$ and $L$ contains a monic polynomial of degree $d-1$;
we recurse. This gives a binary tree of depth at most $d$ and comaximal
elements $s_1,\dots,s_l$ of $A$ such that $L$ has a monic generator $g_j$
over each $A[1/s_j]$. By Lemma~\ref{lem:unique}, which holds over any
ring, $g_j$ and $g_k$ agree over $A[1/(s_js_k)]$, so they glue to a
generator $g$ of $L$. If $A$ is not connected, $g$ need not be monic: its
degree may vary along the leaves. This is the rank-one case of Horrocks'
theorem over an arbitrary ring, the result of Bass and Murthy~\cite{BM}
(the affine Horrocks theorem in rank one,~\cite[V.2]{Lam}). Compared with
the constructive proof following Nashier--Nichols, which splits along a
formal gcd computation and lifts monic polynomials modulo the radical, the
tree here has one binary split per degree and no lifting step. This
extension is formalised (\texttt{DivisorAffineHorrocks}): at each node of
the tree, the retraction is rebuilt from the certificate and $\mu$ is tested
by the existing finite unit/radical test, and the leaves are glued by the
assembly of the formal-gcd proof.
\end{remark}

\section{Rank one}\label{sec:rankone}

Let $L=(y_1,\dots,y_n)$ and suppose $L$ contains a monic polynomial
$F=\sum_kx_ky_k$ with $x_iy_j\in(F)$ for all $i,j$, say $x_iy_j=Fe_{ij}$.
Then $(x_1,\dots,x_n)(y_1,\dots,y_n)=(F)$: $L$ is an invertible ideal,
certified by $F=\sum_ix_iy_i$ itself, without further coefficients. This is
the setting of \cite[Prop.~IV.5.2]{Lam}. Such a certificate is the same as a
rank-one projector: since $F$ is monic, it is a regular element, and
\[
F e_{ik}y_j=x_iy_ky_j=Fe_{ij}y_k,\qquad F\textstyle\sum_ie_{ii}=\sum_ix_iy_i=F
\]
give $e_{ik}y_j=e_{ij}y_k$ and $\sum_ie_{ii}=1$; moreover
$F^2\sum_je_{ij}e_{jk}=\sum_j(x_iy_j)(x_jy_k)=x_iFy_k=F^2e_{ik}$, so
$(e_{ij})$ is idempotent. Conversely, if $e_{ik}y_j=e_{ij}y_k$,
$\sum_ie_{ii}=1$ and $F=\sum_ky_kc_k$, then $x_i=\sum_ke_{ik}c_k$ satisfies
\[
x_iy_j=\textstyle\sum_k e_{ik}c_ky_j=\sum_k e_{ij}y_kc_k = Fe_{ij},
\qquad \sum_iy_ix_i=F\sum_ie_{ii}=F .
\]

\begin{lemma}[splitting]\label{lem:split}
Let $e_{ik}y_j=e_{ij}y_k$, $\sum_ie_{ii}=1$, and $x_iy_j=Fe_{ij}$ with $F$
monic. Then $\rho(p)=\sum_i y_i\,q(x_ip)$, with $q$ the quotient by $F$, is
$A$-linear, takes values in $L$ and is the identity on $L$.
\end{lemma}

\begin{proof}
Linearity comes from that of $q$. If $p=\sum_jy_jb_j$, then
$x_ip=F\sum_je_{ij}b_j$, so $q(x_ip)=\sum_je_{ij}b_j$, and
$\rho(p)=\sum_{i,j}y_ie_{ij}b_j=\sum_{i,j}e_{ii}y_jb_j=p$ using
$y_ie_{ij}=e_{ii}y_j$ and $\sum_ie_{ii}=1$.
\end{proof}

The inverse ideal $L^{-1}$ is generated by the fractions $x_i/F$, and
$q(x_ip)$ is the polynomial part of $(x_i/F)p$, in the decomposition of
fractions with monic denominator into a polynomial and a fraction of
negative degree~\cite[Prop.~IV.1.4]{Lam}. Thus $\rho$ multiplies by
$L^{-1}$, keeps the polynomial part, and multiplies back by $L$.

\begin{corollary}\label{cor:rankone}
Let $y_1,\dots,y_n$, $e_{ij}\in A[X]$ with $e_{ik}y_j=e_{ij}y_k$ and
$\sum_ie_{ii}=1$, and suppose $L=(y_1,\dots,y_n)$ contains a monic
polynomial. Then $L$ is generated by a unique monic polynomial $G\in L$,
that is $y_i=Ga_i$ for some $a_i$.
\end{corollary}

\begin{proof}
Lemma~\ref{lem:split} and Theorem~\ref{thm:main}. The retraction is built
once, from one monic polynomial of $L$, and stays fixed in the induction.
\end{proof}

\begin{corollary}
Let $f\in A[X]$ be monic and $(e_{ij})$ a projector of rank one over $A[X]$
which becomes free after inverting $f$. Then it is free.
\end{corollary}

\begin{proof}
Free after inverting $f$ gives $x_i,y_j$ and $k$ with $x_iy_j=f^ke_{ij}$ and
$\sum_iy_ix_i=f^k$. As above, since $f^k$ is monic, $e_{ik}y_j=e_{ij}y_k$
and $\sum_ie_{ii}=1$, and $L=(y_1,\dots,y_n)$ contains $f^k$. By
Corollary~\ref{cor:rankone}, $y_i=Ga_i$ with $G=\sum_jy_jc_j$. Put
$b_i=\sum_je_{ij}c_j$. Then $x_iG=f^kb_i$, so
$f^kb_ia_j=x_iGa_j=x_iy_j=f^ke_{ij}$ and $f^k\sum_ia_ib_i=\sum_iy_ix_i=f^k$.
Cancelling $f^k$ gives $e_{ij}=b_ia_j$ and $\sum_ia_ib_i=1$: the projector
is free.
\end{proof}

In synthetic algebraic geometry the ring is $R$, which is local, and this
is the form used for chart triviality on $\mathbb P^1$: a line bundle on
$\mathbb A^1$ trivial on $D(f)$, $f$ monic, is trivial.

\begin{remark}
Geometrically, $A[X]/L$ is the coordinate algebra of a closed subscheme of
$\mathbb A^1_A$ contained in $V(F)$, whose coordinate algebra $A[X]/(F)$ is
free of rank $d$. The retraction says that this subscheme is finite and flat
over the base. Lemma~\ref{lem:step} does not change $L$: it lowers the
degree of the containing divisor $V(F)$ until $F$ cuts out the subscheme.
\end{remark}

\section{Ingredients}

Besides polynomial arithmetic, the proof uses only uniqueness of division by
a monic polynomial (also for Lemma~\ref{lem:unique}), the determinant-free NAK
(Lemma~\ref{lem:nak}) and locality of $A$. There is no gluing of local data
and no formal gcd. The splitting $\rho$ was pointed out in a discussion with
ChatGPT.

\begin{sloppypar}
The proof has been checked by machine, in \textsf{Unit}, a small proof
checker for dependent type theory (it is neither Lean nor Agda), designed
by T.~Coquand and implemented by Claude. The formal proof was also written
by Claude. The names in typewriter font are the modules of this
formalisation, in the directory \texttt{sag/horrocks} of
\url{https://github.com/coquand/unit}, in the order of the text: \texttt{MonicReduction} (division), \texttt{UnitNakayama}
(Lemmas~\ref{lem:nak} and~\ref{lem:local}), \texttt{DivisorStep}
(\texttt{Summand}, \texttt{Stable}: Section~\ref{sec:stable}, with locality
only in \texttt{Stable.Local};
\texttt{Retract}: Section~\ref{sec:induction}, with
Theorem~\ref{thm:main} as \texttt{Retract.Induction.generator}),
\texttt{MonicUniqueness} and \texttt{MonicGenerator} (Lemma~\ref{lem:unique};
the type of monic generators is a proposition, \texttt{generator-prop}),
\texttt{DivisorSplitting} (Lemma~\ref{lem:split}) and
\texttt{DivisorHorrocks} (the corollaries). Remark~\ref{rem:dynamic} is
\texttt{DivisorAffineHorrocks}, next to the earlier proof following
Nashier--Nichols (\texttt{HorrocksProof}).
\end{sloppypar}

\begin{thebibliography}{10}
\bibitem{BM} H.~Bass and M.~P.~Murthy, Grothendieck groups and Picard groups
of abelian group rings, \emph{Ann.\ of Math.}\ 86 (1967), 16--73.
\bibitem{Fer} D.~Ferrand, Les modules projectifs de type fini sur un anneau
de polyn\^omes sur un corps sont libres (d'apr\`es D.~Quillen et
A.~A.~Suslin), S\'eminaire Bourbaki 1975/76, Exp.~484, Lecture Notes in
Math.\ 567, Springer, 1977, 202--221.
\bibitem{Hor} G.~Horrocks, Projective modules over an extension of a local
ring, \emph{Proc.\ London Math.\ Soc.}\ (3) 14 (1964), 714--718.
\bibitem{Lam78} T.~Y.~Lam, \emph{Serre's Conjecture}, Lecture Notes in
Math.\ 635, Springer, 1978.
\bibitem{Lam} T.~Y.~Lam, \emph{Serre's Problem on Projective Modules},
Springer, 2006.
\bibitem{Lin} H.~Lindel, Eine Bemerkung zur Quillenschen L\"osung des
Serreschen Problems, \emph{Math.\ Ann.}\ 230 (1977), 97--100.
\bibitem{LQ} H.~Lombardi and C.~Quitt\'e, \emph{Commutative Algebra:
Constructive Methods}, Springer, 2015.
\bibitem{Mat} H.~Matsumura, \emph{Commutative Ring Theory}, Cambridge
University Press, 1986.
\bibitem{NN} B.~Nashier and W.~Nichols, Ideals containing monics,
\emph{Proc.\ Amer.\ Math.\ Soc.}\ 99 (1987), 634--636.
\end{thebibliography}

\end{document}
